% slicing_theory.tex — appendix: the sliced regularizer against the full-covariance one (draft 2026-09-17; theory only,
% per Berker). \input after \section{Implementation of the regularizer} (app:slicing) or as its own section.
% Notation follows the method box (R_ctr; mu, Sigma of the view centers) and app:slicing (Q in R^{d x d'}, d').
% Facts (i)-(iii) are elementary; (iv)-(v) use the operator Jensen inequality for the operator-concave logarithm
% (Hansen & Pedersen 2003). A numeric check on stored features (scratch/slice_check.py, job 66274754) agrees with every
% line: mean and trace parts identical at every d', the log-det part increasing in d' and never above the full value.
\section{Slicing and the full-covariance regularizer}
\label{app:slicing-theory}

Let $\mu\in\mathbb R^d$ and $\Sigma\succ0$ be the mean and covariance of the view centers, and write the full regularizer per dimension as
\[
\mathcal R(\mu,\Sigma)
=\frac{1}{2d}\Big[\|\mu\|_2^2+\operatorname{tr}\Sigma-d-\log\det\Sigma\Big]
=\frac{1}{d}\,\mathrm{KL}\big(\mathcal N(\mu,\Sigma)\,\|\,\mathcal N(0,I_d)\big).
\]
For an orthonormal $Q\in\mathbb R^{d\times d'}$ the projected centers have mean $Q^\top\mu$ and covariance $Q^\top\Sigma Q$, so the sliced regularizer of Appendix~\ref{app:slicing} is the same divergence for the projected Gaussian,
\[
\mathcal R_Q(\mu,\Sigma)
=\frac{1}{2d'}\Big[\|Q^\top\mu\|_2^2+\operatorname{tr}(Q^\top\Sigma Q)-d'-\log\det(Q^\top\Sigma Q)\Big]
=\frac{1}{d'}\,\mathrm{KL}\big(\mathcal N(Q^\top\mu,Q^\top\Sigma Q)\,\|\,\mathcal N(0,I_{d'})\big).
\]
Since $Q$ is redrawn at every step from the uniform (Haar) distribution on orthonormal frames, training minimizes the expected sliced regularizer
\[
\bar{\mathcal R}_{d'}(\mu,\Sigma)=\mathbb E_Q\big[\mathcal R_Q(\mu,\Sigma)\big]
\]
by stochastic gradient descent. The following statement relates $\bar{\mathcal R}_{d'}$ to $\mathcal R$.

\begin{proposition}[Slicing]
\label{prop:slicing}
For every $1\le d'\le d$:
\begin{enumerate}[leftmargin=*]
\item[(i)] \emph{Same minimizer.} $\bar{\mathcal R}_{d'}\ge0$, with equality if and only if $\mu=0$ and $\Sigma=I_d$.
\item[(ii)] \emph{Unbiased gradient.} If $\mu$ and $\Sigma$ depend smoothly on parameters $\theta$, then $\mathbb E_Q\nabla_\theta\mathcal R_Q=\nabla_\theta\bar{\mathcal R}_{d'}$: a fresh slice per step is an unbiased stochastic gradient of $\bar{\mathcal R}_{d'}$.
\item[(iii)] \emph{Mean and scale.} $\mathbb E_Q\|Q^\top\mu\|_2^2/d'=\|\mu\|_2^2/d$ and $\mathbb E_Q\operatorname{tr}(Q^\top\Sigma Q)/d'=\operatorname{tr}\Sigma/d$: the mean and trace terms of the sliced regularizer are unbiased for those of $\mathcal R$.
\item[(iv)] \emph{Log-determinant.} $\frac{1}{d'}\mathbb E_Q\log\det(Q^\top\Sigma Q)\ge\frac{1}{d}\log\det\Sigma$; consequently $\bar{\mathcal R}_{d'}\le\mathcal R$, with equality at $d'=d$, where $\mathcal R_Q=\mathcal R$ for every $Q$.
\item[(v)] \emph{Monotonicity.} $\bar{\mathcal R}_{d'}$ is nondecreasing in $d'$, so $\bar{\mathcal R}_{d'}\uparrow\mathcal R$ as $d'\to d$.
\end{enumerate}
\end{proposition}

\begin{proof}[Proof sketch]
(i) Each $\mathcal R_Q$ is a KL divergence, hence nonnegative, and vanishes if and only if $Q^\top\mu=0$ and $Q^\top\Sigma Q=I_{d'}$. If $\bar{\mathcal R}_{d'}=0$ these hold for almost every frame and, by continuity in $Q$, for every frame; a frame whose first column is an arbitrary unit vector $u$ gives $u^\top\mu=0$ and $u^\top\Sigma u=1$ for all $u$, i.e.\ $\mu=0$ and $\Sigma=I_d$. (ii) The set of frames is compact and $\mathcal R_Q$ is smooth in $(\mu,\Sigma)$ on $\Sigma\succ0$ (the $\varepsilon I$ of the implementation keeps this on the closure), so differentiation and expectation commute. (iii) Haar invariance implies that $\mathbb E[QQ^\top]$ commutes with every orthogonal matrix and has trace $d'$, hence $\mathbb E[QQ^\top]=\frac{d'}{d}I_d$; both identities follow by linearity. (iv) Write $\Sigma=\sum_i\lambda_ie_ie_i^\top$ and $v_i=Q^\top e_i$, so that $Q^\top\Sigma Q=\sum_i\lambda_iv_iv_i^\top$ with $\sum_iv_iv_i^\top=I_{d'}$. The operator Jensen inequality for the operator-concave logarithm gives $\log(Q^\top\Sigma Q)\succeq\sum_i(\log\lambda_i)\,v_iv_i^\top$; taking the trace, $\log\det(Q^\top\Sigma Q)\ge\sum_i\|v_i\|^2\log\lambda_i$, and $\mathbb E\|v_i\|^2=d'/d$ by (iii). Together with (iii) this gives $\bar{\mathcal R}_{d'}\le\mathcal R$. For $d'=d$, $Q$ is orthogonal, $Q^\top\Sigma Q$ has the eigenvalues of $\Sigma$ and $\|Q^\top\mu\|=\|\mu\|$. (v) A uniformly random $d'$-frame is a uniformly random sub-frame of a uniformly random $(d'+1)$-frame; applying (iv) inside that $(d'+1)$-dimensional slice gives $\frac{1}{d'}\mathbb E\log\det(\text{$d'$-slice})\ge\frac{1}{d'+1}\mathbb E\log\det(\text{$(d'+1)$-slice})$, while the mean and trace terms do not depend on $d'$ by (iii).
\end{proof}

In words: at every slice width the regularizer has the same unique minimizer as the full one, the isotropic unit-covariance Gaussian, and training follows unbiased stochastic gradients of a well-defined objective. The mean and the overall scale are seen exactly by every slice; only the log-determinant, the part that penalizes vanishing directions, is seen partially, more fully as $d'$ grows and completely at $d'=d$. Within each slice the divergence is two-sided along all of its $d'$ directions, so at any width every projected variance is pushed toward one from above and from below. The slice width therefore trades how much of the spectrum is penalized per step against the number of centers needed to estimate the projected covariance, which is the reason for the choice $d'<d$ in Appendix~\ref{app:slicing}. This is the same construction as sliced divergences between distributions: for Gaussians, the projections onto all frames determine the distribution, which is (i) in another form.
